
\iffalse
\begin{table}[t]
\centering
\small
\caption{Time complexity of various set operations. $n$ is the number of items, $d$ is the dimensionality of each item, and $m$ is the number of inducing points.}
\label{tab:complexities}
\begin{center}
\vspace{-1mm}

\begin{tabular}{lccc}
\toprule
Set operations & Time complexity & High-order & Permutation \\
&& interactions & invariant \\
\hline
%Fully- Connected (naive) & $O(n^2d^2)$ & Yes & No \\
\rule{0pt}{2.0ex}Recurrent & $O(nd)$ & Yes & No \\
Pooling~\citep{Zaheer2017} & $O(nd)$ & No & Yes \\
Relational Networks~\citep{Santoro2017} & $O(n^2d)$ & Yes & Yes \\
\midrule
Set Transformer (SAB + PMA, ours) & $O(n^2d)$ & Yes & Yes \\
Set Transformer (ISAB + PMA, ours) & $O(nmd)$ & Yes & Yes \\
\bottomrule
\end{tabular}
\end{center}
\end{table}
\fi


\if(0)
We first introduce the \textit{set attention block} (SAB), a basic building block of the Set Transformer.  Unlike the element-wise encoding methods for sets, the set transformer encodes entire sets altogether with using self-attention mechanism. More specifically, given matrices $X, Y \in \R^{n\times d}$ collecting the vector representation of elements in two sets, SAB is defined as
\[
\mathrm{SAB}(X, Y) &= \mathrm{LayerNorm}(H + \mathrm{FF}(H)), \\
\textrm{where } H &= \mathrm{LayerNorm}(X + \mathrm{Multihead}(X, Y, Y)),
\]
where $\mathrm{FF}$ is any type of feed-forward layer that acts on each element identically and independently, and $\mathrm{LayerNorm}$ is layer normalization~\citep{Ba2016}. The SAB is a simple adaptation of the building blocks of transformer architecture without positional encoding and dropout.  We use the simplified notation $\mathrm{SAB}(X) \coloneqq \mathrm{SAB}(X, X)$ for self-attention using the $\mathrm{SAB}$.

A simple encoder structure is a stack of SABs:
\[
\mathrm{Encoder}(X) = \mathrm{SAB}(\mathrm{SAB}(\cdots \mathrm{SAB}(X))).
\]
By stacking SABs, we can encode high-order interactions between the elements in the set $X$ (recall that the matrix $X$ represents a set of size $n$).
A problem with using SABs for set data is the time complexity $O(n^2)$, quickly making this method intractable for large sets ($n \gg 1$).

We thus introduce the \textit{induced set attention block} (ISAB), which is an attention-based encoding block that can bypass this problem.
Let $X\in \R^{n \times d}$ be our data and define $m$ new trainable $d$-dimensional vectors ($I \in \R^{m \times d}$).
An ISAB is defined as:
\[
\mathrm{ISAB}(X;I) \coloneqq \mathrm{SAB}(I, X).
\]
Note that attention was computed between a set of size $m$ and a set of size $n$.
Therefore, while the time complexity of self-attention $\mathrm{SAB}(X)$ is $O(n^2)$, the time complexity of $\mathrm{ISAB}(X;I)$ is $O(nm)$, where $m$ is a constant hyperparameter.
We compare characteristics of various set operations in Table \ref{tab:complexities}.

The idea of introducing trainable vectors $I$ is loosely based on the inducing point methods used in sparse Gaussian process regression~\citep{Quinonero-Candela2005} and the Nystr\"{o}m method for matrix decomposition~\citep{Fowlkes2004}.


\subsection{Encoder}
\label{subsec:imab}
We construct the encoder of the Set Transformer by alternating $\mathrm{IMAB}$ and $\mathrm{MAB}$ blocks, for example:
\begin{align}
\mathrm{Encoder}(X) &= \mathrm{MAB}(X, H) \in \R^{n\times d} \\
\mathrm{where} \hspace{10pt} H &= \mathrm{IMAB}(X;I)\in \R^{m\times d},
\end{align}
where $I \in \R^{m \times d}$.
Note that $H \in \R^{m\times d}$, so the $\mathrm{MAB}$ operation has time complexity $O(nm)$.

\subsection{Decoder}
After the encoder transforms data $X \in \R^{n \times d}$ into feature representations $Z \in \R^{n\times d}$,
the decoder then aggregates these features into a single vector to fed into a feed-forward network to get final outputs.
As mentioned in section \ref{subsec:pooling}, a commonly used method of aggregating these features is to take sums or averages of the feature vectors.
The set transformer's decoder uses an attention-based pooling mechanism using the MAB.
We prove in Appendix \ref{app:proofs} that commonly used pooling layers such as mean-pooling and max-pooling are special cases of our decoder structure.
Similarly to how recurrent neural nets take an initial hidden state and decode it into an output, we use an IMAB with one inducing point $s \in \R^d$:
\[
\mathrm{Decoder}(X) &= \mathrm{FF}(H) \\
\mathrm{where} \hspace{10pt} H &= \mathrm{IMAB}(X; s).
\]
The IMAB here can be understood as a pooling operation using the weighted sum of features from each head.

For some problems such as clustering, we want to decode $k$ correlated outputs.
In this case, we introduce a set of seed vectors $S \in \R^{k \times d}$, and apply our blocks as follows:
\[
\mathrm{Decoder}(X) &= \mathrm{FF}(H) \in \R^{k \times d_{\mathrm{out}}} \\
\textrm{ where } \hspace{10pt} H &= \mathrm{MAB}(\mathrm{IMAB}(Z; S)) \in \R^{k \times d}.
\]
The role of the MAB here is to enable the decoder to model high-order interactions between its $k$ outputs.
\fi



\iffalse
Unlike previous work, we explicitly parameterize pairwise interactions in each self-attention block.
Because our proposed model consists of several such blocks, it can encode high-order interactions among elements in a set.
Additionally, we introduce the ISAB, which reduces the $O(n^2)$ time complexity of full self-attention to $O(n)$.
\fi

\iffalse
\cite{Vinyals2016a} proposes an architecture to map sets into sequences,
where elements in a set are pooled by weighted average with weights computed from attention mechanism.
Several recent works have highlighted the competency of attention mechanisms in modeling sets.
\citep{Yang2018} proposes AttSets for multi-view 3D reconstruction, where attention is applied to the encoded features of elements in sets before pooling.
Similarly, \citep{Ilse2018} uses an attention in pooling for multiple instance learning.
However, the attention applied in those papers are simple dot-product attention that does not encode higher-order interactions between elements.
Although not permutation invariant, \citep{Mishra2018} has an attention as one of its core components to meta-learn to solve various tasks using sequences of inputs.
\fi





